\chapter{Background}
\label{ch:background}

This chapter starts off with \autoref{sec:network-protocol-background}, which introduces
two fundamental concepts: the byte order and the \code{VarInt} data type.
Following this, a brief introduction of the Bitcoin network, necessary to understand the \tech{ln-history} platform, is given in \autoref{sec:bitcoin-network}.
After elaborating on the scalability limitations of the Bitcoin blockchain,
we motivate in \autoref{sec:lightning-network} the development of the \gls{ln} as potential scaling solution.
The Lightning Protocol is defined in the Basis of Lightning Technology (BOLT)~\cite{lightningdevelopers_bolt-3_2026}.
In \autoref{sec:lightning-protocol}, we present the overall structure of the \gls{ln} and 
focus on BOLT\#7, the part of the specification that defines the gossip protocol.
In essence three different messages types: 
\code{channel\_announcement}, \code{node\_announcement} and \code{channel\_update}
form the foundational data for the \tech{ln-history} platform.
As the specifications need to be implemented, we present the current existing client 
implementations at the end in \autoref{sec:implementations}.


\section{Network Protocol Background}
\label{sec:network-protocol-background}

\subsection{Byte Order}
When transmitting data across a network, systems must agree on a specific \term{byte order}. 
As defined by \citeauthor{patterson_computer_2014}~\cite{patterson_computer_2014}, 
this is the convention a machine uses to number bytes within a word. 
In \term{\gls{be}} format, the most significant byte is transmitted first, 
whereas \term{\gls{le}} format transmits the least significant byte first.
This distinction frequently creates friction between hardware and network engineering. 
While modern processor architectures (such as x86 and ARM) predominantly execute in \gls{le}, 
standard network communication protocols enforce \gls{be}~\cite{reynolds_rfc_1994}. 

The Bitcoin base layer also has this architectural inconsistency: 
its core data structures were originally optimized by \citeauthor{nakamoto_bitcoin_2008} for x86 (\gls{le}), 
yet user-facing cryptographic identifiers, such as \code{txids}, 
are displayed in \gls{be} to preserve the human readability of leading zeros~\cite{antonopoulos_bitcoinbook_2024}. 
Consequently, to parse network topology accurately, any Lightning node or data ingestion platform must 
manage endianness depending on the exact protocol layer and message field being evaluated.


\subsection{Variable Integer (\code{VarInt})}
The data type used for integers in Bitcoin is \gls{varint} which does not use a 
fixed number of bits to store a value but dynamically adapts its length depending 
on the order of magnitude of the number. 
As detailed in ~\autoref{tab:varint_rules} values smaller than \code{0xFD} are stored in a single byte. 
Larger values need a specific one-byte prefix (one of \code{0xFD}, \code{0xFE}, or \code{0xFF}) 
depending on their value~\cite{antonopoulos_bitcoinbook_2024}.

\begin{table}[htbp]
    \centering
    \caption{Encoding rules for Bitcoin Variable Length Integers.}
    \label{tab:varint_rules}
    \begin{tabular}{|l|l|l|}
        \hline
        \textbf{Value Range ($N$)} & \textbf{Prefix (1 Byte)} & \textbf{Encoded Payload} \\ \hline
        $0 \le N \le \code{0xFC}$ & \textit{None} & $N$ (as uint8) \\ \hline
        $\code{0xFD} \le N \le \code{0xFFFF}$ & \code{0xFD} & $N$ as \gls{le} uint16 \\ \hline
        $\code{0x10000} \le N \le \code{0xFFFFFFFF}$ & \code{0xFE} & $N$ as \gls{le} uint32 \\ \hline
        $N \ge \code{0x100000000}$ & \code{0xFF} & $N$ as \gls{le} uint64 \\ \hline
    \end{tabular}
\end{table}

It is important to note that the \code{VarInt} standard in Bitcoin does not support 
negative numbers or integers exceeding the 64-bit unsigned limit ($2^{64}$).
Further details can be found directly in the codebase of the Bitcoin Core software at 
\code{src/serialize.h}~\cite{bitcoincoredevelopers_bitcoin_2026}. 


\section{Bitcoin}
\label{sec:bitcoin-network}

On 31st October 2008 the pseudonym Satoshi Nakamoto published the Bitcoin whitepaper 
\citetitle{nakamoto_bitcoin_2008}~\cite{nakamoto_bitcoin_2008}.
\citeauthor{nakamoto_bitcoin_2008} combined several existing concepts into a novel idea: 
A digital \gls{p2p} system for trustless, borderless, and pseudonymous
transactions using an absolutely scarce unit of account.

Transactions form the foundational layer of the protocol, enabling the secure transfer of value. 
To prevent the double-spending problem without relying on a centralized entity, 
transactions are batched into blocks. 
These blocks are cryptographically chained together via a Proof-of-Work consensus mechanism, 
creating a chronological, immutable ledger of events where every 
unit of value can be mathematically traced back to its origin.

\subsection{Transaction Model}
\label{sec:transaction}

\autoref{fig:bitcoin_trasaction_visual} visualizes the concept of a bitcoin transaction.
Unlike traditional banking systems that use account balances, 
Bitcoin utilizes a \gls{utxo} model.

\begin{figure}[htpb]
    \centering
    \tikzset{
        box/.style={draw, minimum size=2em, text width=3em, text centered},
        container/.style={draw, inner sep=20pt}
    }

    \begin{tikzpicture}[>=Triangle]
    \foreach[evaluate={\prev=int(\i - 1)}] \i in {1, 2, 3} {
        \node (P\i-PubKey) at (\i * 5, 0) [box, text width=5em] {Owner \i's \\ Public Key};
        \node (P\i-Hash) [below=0.8cm of P\i-PubKey][box] {Hash};
        \node (P\i-Sig) [below=of P\i-Hash][box, text width=5em] {Owner \prev's \\ Signature};
        \draw [->] (P\i-Hash) -- (P\i-Sig);
        \draw [->] (P\i-Hash.45 |- P\i-PubKey.south) -- (P\i-Hash.45);
        \node (P\i-Tx) [fit=(P\i-PubKey)(P\i-Hash)(P\i-Sig)][container, label={[shift={(15ex,-4ex)}]north west:Transaction}]{};
        \node (P\i-PrivKey) [below=of P\i-Tx] [box, text width=6em] {Owner \i's \\ Private Key};
    }

    \draw [->] ([xshift=-2.5cm, yshift=0.5cm]P1-Hash.north) -| ([xshift=-0.3cm]P1-Hash.north);

    \foreach[evaluate={\next=int(\i + 1)}] \i in {1, 2} {
        \draw [->, dashed] (P\i-PubKey.330) -- (P\i-PubKey.330 |- P\i-Hash.south east) -- node[sloped]{Verify} (P\next-Sig.165);
        \draw [->, dashed] (P\i-PrivKey.15) -- node[sloped]{Sign} (P\next-Sig.195);
        \draw [->] ([yshift=0.5cm]P\i-Hash.north -| P\i-Tx.east) -| ([xshift=-0.3cm]P\next-Hash.north);
    }
    \end{tikzpicture}
    \caption{
        Chain of ownership structure in Bitcoin.
        Each owner transfers the coin to the next by digitally signing a hash of 
        the previous transaction and the public key of the next owner~\cite{nakamoto_bitcoin_2008}.
    }
    \label{fig:bitcoin_trasaction_visual}
\end{figure}


In essence, a transaction consists of metadata alongside two lists: 
a list of \term{inputs} and a list of \term{outputs}. 
Every broadcasted transaction is uniquely identifiable by its cryptographic hash, 
known as the \code{txid} (transaction ID). 
An output has a specific amount of \term{satoshis} and assigns 
a cryptographic locking script (the \code{scriptPubKey})
determining the conditions required to spend it. 
The global state of the Bitcoin network is the set of all currently existing \glspl{utxo}.
To spend a \gls{utxo}, a user constructs a new transaction where the input provides 
the exact \code{txid} and the specific index (\code{vout}) of the previous output, 
alongside an unlocking script (a digital signature) to satisfy the cryptographic lock.


\subsection{Timechain and Consensus}
\label{sec:timechain}

To maintain global consensus, unconfirmed transactions broadcasted to the \gls{p2p} network 
are collected by miners and batched into discrete data structures called blocks. 
As illustrated in \autoref{fig:timechain}, this architecture functions as a distributed timestamp server, 
where each block includes the cryptographic hash of the previous block's timestamp, 
creating a chronological and immutable chain.

\begin{figure}[htpb]
    \centering
    \begin{tikzpicture}[>=Triangle]
    \tikzset{box/.style={draw, minimum size=2em, text width=2em, text centered},
        container/.style={draw, inner sep=20pt}
    }

    \foreach \i / \x in {0/-2, 1/4} {
        \node (P\i-Hash) [box] at (\x, 1.5) {Hash};
        \node (Item1) [box] at (\x - 1, 0){Item};
        \node (Item2) [box] [right=0.3cm of Item1] {Item};
        \node (Item3) [box] [right=0.3cm of Item2] {...};
        \node (Container) [container] [label={[shift={(8ex,-4ex)}]north west:Block}, fit=(Item1)(Item2)(Item3)] {};

        \draw [->] (Container.north west)+(1, 0) |- ([yshift=-4]P\i-Hash.west);
    }

    \draw [<-] ([yshift=4]P0-Hash.west) --+ (-3, 0);
    \draw [->] ([yshift=4]P0-Hash.east) -- ([yshift=4]P1-Hash.west);
    \draw [->] ([yshift=4]P1-Hash.east) --+ (3, 0);
    \end{tikzpicture}
    \caption{
        The abstract concept of the distributed timestamp server. 
        Diagram recreated based on \citeauthor{nakamoto_bitcoin_2008}~\cite{nakamoto_bitcoin_2008}.
    }
    \label{fig:timechain}
\end{figure}

To secure this timechain against tampering, Bitcoin employs a Proof-of-Work consensus mechanism. 
As detailed in \autoref{fig:proof-of-work}, miners must expend computational energy in the physical world 
to find a valid \code{hash} that satisfies the network's difficulty target. 
Once a valid block is mined and appended to this chain of blocks, hence blockchain, the transactions within it are considered 
confirmed\footnote{Because blockchain reorganizations (\term{reorgs}) can occasionally occur due to network latency, 
\citeauthor{nakamoto_bitcoin_2008} mathematically showed that waiting for a transaction to be six blocks 
deep provides probabilistic finality~\cite[Section~11]{nakamoto_bitcoin_2008}.}.

\begin{figure}[htpb]
    \centering
    \begin{tikzpicture}[>=Triangle]
    \tikzset{box/.style={draw, minimum size=2em, text centered},
        container/.style={inner sep=20pt}
    }

    \foreach \i / \x in {0/-3, 1/3} {
        \node (P\i-PrevHash) [box] at (\x, 1) {Prev Hash};
        \node (Nonce) [box] [right=0.3cm of P\i-PrevHash] {Nonce};
        \node (Tx1) [box] [below=0.3cm of P\i-PrevHash]{Tx};
        \node (Tx2) [box] [right=0.3cm of Tx1] {Tx};
        \node (Tx3) [box] [right=0.3cm of Tx2] {...};
        \node (P\i-Container) [container] [label={[shift={(8ex,-4ex)}]north west:Block}, fit=(P\i-PrevHash)(Tx1)(Tx2)(Tx3)] {};
        \draw (P\i-Container.north west) -- (P\i-Container.south west) -- (P\i-Container.south east) -- (P\i-Container.north east);
    }

    \draw [<-] (P0-PrevHash.west) --+ (-2, 0);
    \draw [->] (P0-Container.east |- P1-PrevHash.west) -- (P1-PrevHash.west);
    \end{tikzpicture}
    \caption{
        The Proof-of-Work blockchain architecture, illustrating how transaction batches are cryptographically chained using the previous block's hash and a valid Nonce. Diagram recreated based on Nakamoto~\cite{nakamoto_bitcoin_2008}.
    }
    \label{fig:proof-of-work}
\end{figure}

Since these blocks are strictly ordered and the transactions within a block are sequentially listed, 
any confirmed output in the Bitcoin network can be deterministically located. 
This forms a three-dimensional tuple: 
the \term{block height}, the \term{transaction index} within that block, and the \term{output index} (\code{vout}) within that transaction. 
As will be detailed in later chapters, the \gls{ln} leverages this tuple,
encoded as a \code{short\_channel\_id}, to concisely identify and verify the existence of payment channels on the base layer.


\section{Scalability of Bitcoin}

By design, the Bitcoin protocol limits the size of each block to ensure that running 
a fully validating node remains computationally accessible to the public, thereby preserving network decentralization. 
Originally, a strict \qty{1}{\mega\byte} storage limit was enforced. 
As global adoption increased, this hard limit created a severe bottleneck in transaction throughput, 
leading to highly congested mempools and surging transaction fees.

These scalability challenges triggered a conflict known as the \enquote{Blocksize War}~\cite{bier_blocksize_2021}. 
A group of corporate entities and mining pools advocated for a \term{hard-fork} to increase the block size, 
while independent researchers and protocol developers argued for maintaining the limit and scaling via layered architectural upgrades. 
The developers' approach ultimately succeeded, resulting in the 2017 activation of 
a backwards-compatible soft-fork known as Segregated Witness (\term{SegWit}).

SegWit fundamentally reorganized transaction data by separating
the cryptographic signatures (the witness) from the transaction inputs. 
Most importantly for Layer-2 scalability, SegWit eliminated third-party \term{transaction malleability},
a vulnerability where an attacker could alter a \code{txid} before it was confirmed. 
By fixing this, SegWit finally made it secure to chain unconfirmed transactions together.
Furthermore, the protocol relies on \glspl{htlc}, 
which are specialized Bitcoin scripts that lock funds using hash preimages and time expirations. 
Together with the activation of SegWit, these changes laid the 
foundation for a scalable, off-chain payment network: the Lightning Network.


\section{Bitcoin Lightning Network}
\label{sec:lightning-network}

The Bitcoin Lightning Network (LN) was first proposed in 2016 by \citeauthor{poon_bitcoin_2016} 
in their whitepaper, \citetitle{poon_bitcoin_2016}~\cite{poon_bitcoin_2016}. 
Designed as a Layer-2 scaling solution, the \gls{ln} mitigates the throughput limitations
of the Bitcoin blockchain by shifting the majority of transaction volume \term{off-chain}. 
This architecture significantly conserves scarce block space while continuing
to inherit the robust, decentralized security model of the underlying network.

The fundamental building blocks of this off-chain network are \term{payment channels}. 
Technically, a payment channel is a 2-of-2 multisignature transaction, referred to as \term{funding transaction}, 
broadcasted to the Bitcoin base layer. 
Once this channel is established, the two participating nodes can exchange transactions 
by continuously signing and exchanging new \term{commitment transactions} that update the distribution of the locked funds. 
Because of SegWit, these unconfirmed state updates can be exchanged securely 
and trustlessly without requiring immediate on-chain settlement.

To make individual payment channels function as a globally routable payment network, 
nodes must follow a set of communication and cryptographic standards. 
In the following section, we elaborate on the protocol specifications, known as \glspl{bolt}. 
We specifically focus on the gossip protocol that forms the data foundation for the \tech{ln-history} platform.
We finish with an overview of the client software that implement these standards.

\section{Lightning Protocol: BOLT}
\label{sec:lightning-protocol}

The technical specifications of the \gls{ln} are formalized as the BOLTs~\cite{lightningdevelopers_lightning_2018}. 
Currently, comprising twelve documents BOLT\#1 through BOLT\#12, 
these standards ensure interoperability across different client implementations. 
Because the \tech{ln-history} platform is fundamentally designed to parse and store network topology data, 
we first introduce the protocol's core data primitives (\autoref{sec:lightning-data-model}) 
before detailing the specific gossip messages used for network discovery.


\subsection{Data Model}
\label{sec:lightning-data-model}

\subsubsection{Type-Length-Value Records (TLV)}
To ensure backwards compatibility and allow for continuous protocol upgrades, 
Lightning messages extensively utilize \gls{tlv} streams. 
Rather than relying on immutable data structures, messages append optional fields using the format illustrated in \autoref{fig:tlv_record}.

\begin{figure}[htbp]
    \centering
    \renewcommand{\arraystretch}{1.5} 
    
    \begin{tabular}{|c|c|c|}
        \hline
        % Row 1: The Field Names
        \textbf{Type} & \textbf{Length} & \textbf{Value} \\
        
        % Row 2: The Encoding / Size
        \small \textit{\gls{varint}} & \small \textit{\gls{varint}} & \small \textit{Variable ($N$ bytes)} \\
        \hline
    \end{tabular}
    \caption{
        Structure of a single TLV Record. 
        The \term{Value} field's size is determined by the integer in the \term{Length} field.
    }
    \label{fig:tlv_record}
\end{figure}

A \gls{tlv} record consists of a \term{Type} identifier and a \term{Length} indicator, 
both encoded as \gls{varint}, followed by the actual \term{Value} payload. 
Types below $2^{16}$ are reserved for official BOLT specifications, 
while higher values are available for custom or experimental application data~\cite{lightningdevelopers_lightning_2018}. 
By concatenating multiple records into a stream, nodes can safely parse known fields and 
safely ignore unrecognized types without breaking message decoding.


\subsubsection{Short Channel ID (\code{scid})}

As introduced in \autoref{sec:timechain}, 
payment channels are anchored to the Bitcoin base layer via a funding transaction. 
The \gls{scid} provides a compact, unique identifier for these channels 
by encoding their exact on-chain location into an \qty{8}{\byte} integer~\cite{lightningdevelopers_lightning_2018}.

The structure consists of the \term{block height} (\qty{3}{\byte}), 
the \term{transaction index} within that block (\qty{3}{\byte}), 
and the \term{output index} (\qty{2}{\byte}). 
For human readability, these three components are printed as decimal numbers separated by an \code{x}. 
For example, a \gls{scid} of \code{539268x845x1} indicates a channel funded in block \num{539268}, 
at transaction index \num{845}, utilizing output \num{1}\footnote{Early versions of the \gls{cln} 
implementation used a colon (\code{:}) as a separator, but this was standardized to \code{x} 
within the formal BOLT specifications~\cite{user1692333_how_2022}.}.


\subsection{Messaging: BOLT\#1}
\label{sec:lightning-message}

After introducing the core data primitives, we now examine how nodes actually communicate. 
The foundational specification is BOLT\#1 that outlines the base protocol 
for connection and data exchange within the \gls{ln}. 
Because all communication over the network is strictly encrypted and authenticated, 
the protocol formally defines a plaintext \term{Lightning Message Format} that 
emerges only after a packet is decrypted~\cite{lightningdevelopers_bolts_2018b}. 

\begin{figure}[htbp]
    \centering
    \begin{bytefield}[bitwidth=1.1em, endianness=big]{16}
        
        \bitheader{0, 15} \\
        
        %  Type (2 bytes = 16 bits) 
        \bitbox{16}{Type (2 Bytes)} \\
        
        %  Payload (Variable Length) 
        \wordbox[lrt]{1}{Payload (Variable Length)} \\
        \wordbox[lrb]{1}{\footnotesize \textit{Conforms to format matching the Type}} \\
        
        %  Extension (Variable Length / Optional) 
        \wordbox{2}{Extension (Optional TLV Stream)}
        
    \end{bytefield}
    \caption{Structure of a Lightning Message. \cite[section ~Lightning Message Format]{lightningdevelopers_lightning_2018}}
    \label{fig:lightning_message}
\end{figure}


As shown in \autoref{fig:lightning_message}, every parsed message follows to a hierarchical structure. 
The message always begins with a \qty{2}{\byte} \term{Type} field (encoded as a \gls{varint} integer),
which acts as an identifier to dictate how the subsequent variable-length \term{Payload} must be parsed. 
Finally, to ensure future protocol extensibility, messages may optionally conclude with a \gls{tlv} stream.



\subsection{P2P Node and Channel Discovery: BOLT\#7}
\label{sec:gossip-messages}

The discovery of nodes and channels which form the topology of the \gls{ln} does not rely on any third party. 
Instead, nodes use \term{(Public) Gossip Messages} which are \term{Lightning Messages} 
of type \code{256}, \code{257} or \code{258} and defined in ~\cite{lightningdevelopers_bolts_2018a}. 
The gossip protocol serves two objectives

\begin{enumerate}
    \item \textbf{Peer discovery:} 
        Every node has the properties \code{node\_id}, \code{alias}, \code{color}, 
        and a list of \code{addresses} that other nodes can use to connect to.
    \item \textbf{Channel discovery:} 
        Every node creates and maintains its own (local) view 
        of the topology of the \gls{ln}, necessary for path finding and routing.
\end{enumerate}

To fulfil these objectives, the protocol defines three distinct message types:

Firstly, the \code{channel\_announcement} (type \code{256}) establishes the topological link between two peers. 
It binds the participating nodes (\code{node\_id\_1} and \code{node\_id\_2}) to a specific on-chain \gls{scid}. 
To prevent malicious actors from spamming the network with fake topology data, 
this message requires four distinct cryptographic signatures (from both the node keys and the underlying Bitcoin funding keys). 
This guarantees that the announced channel is genuinely anchored to a confirmed Bitcoin transaction.

Secondly, the \code{node\_announcement} (type \code{257}) is broadcasted by individual nodes to declare their network reachability. 
It provides the necessary routing \code{addresses} alongside the node's metadata. 
To mitigate spam, the protocol mandates that a node may only broadcast this message 
after it has successfully verified and announced at least one public channel.

Thirdly, the \code{channel\_update} (type \code{258}) informs the network about 
the operational and economic policies of an established channel. 
Because \gls{ln} channels are bidirectional, each participating node broadcasts its own update to 
independently set the routing fees (base and proportional) and timelock requirements (\term{HTLC} minimums and maximums) 
for its specific forwarding direction. 
These updates make up the vast majority of network traffic as nodes continuously adjust their fee policies to manage their liquidity.


\subsubsection{Channel Announcement}
\label{sec:channel-announcement}

The \code{channel\_announcement} message broadcasts the existence of 
a newly established, public channel to the \gls{ln}. 
It links the founding Lightning nodes to the on-chain funding transaction that backs their channel. 
The structure of the message payload is defined by the BOLT specification and is illustrated in \autoref{fig:channel_announcement}.

\begin{figure}[htbp]
    \centering
    % 32-bit width
    \begin{bytefield}[bitwidth=1.1em, endianness=big]{32}
        \bitheader{0, 15, 16, 31} \\
        
        % Type Header 
        \wordbox{1}{\code{type} (2 bytes, value: 256)} \\

        % Signatures 
        \wordbox[lrt]{1}{4 $\times$ 64-byte signatures} \\
        \wordbox[lrb]{1}{\code{node\_sig\_1}, \code{node\_sig\_2}, \code{bitcoin\_sig\_1}, \code{bitcoin\_sig\_2}} \\
        
        % Features 
        \bitbox{16}{\code{len} (2 bytes)} & \bitbox{16}{\code{features} (\code{len} bytes)} \\
        
        % Chain Hash 
        \wordbox{1}{\code{chain\_hash} (32 bytes)} \\
        
        % SCID 
        \wordbox{2}{\code{short\_channel\_id} (8 bytes)} \\

        % Node IDs 
        \wordbox{1}{\code{node\_id\_1} (33 bytes)} \\
        \wordbox{1}{\code{node\_id\_2} (33 bytes)} \\
        
        % Bitcoin Keys
        \wordbox{1}{\code{bitcoin\_key\_1} (33 bytes)} \\
        \wordbox{1}{\code{bitcoin\_key\_2} (33 bytes)}
        
    \end{bytefield}
    \caption{
        Byte layout of the \code{channel\_announcement} message. 
        Large fields are shown as blocks~\cite[Section~channel\_announcement]{lightningdevelopers_lightning_2018}.
    }
    \label{fig:channel_announcement}
\end{figure}

The fields are defined as follows:
\begin{description}
    \item[Signatures] 
        Four \qty{64}{\byte} cryptographic signatures. 
        Together, they prove that both Lightning nodes agree on the channel and 
        that they own the associated on-chain funding outputs. 
    
    \item[\code{features}] 
        A \gls{varint} bitfield used to signal supported protocol extensions (defined in BOLT\#9~\cite{lightningdevelopers_bolts_2018}). 
        This ensures forward compatibility, allowing older nodes to process announcements from newer nodes without breaking network consensus.

    \item[\code{chain\_hash}] 
        A \qty{32}{\byte} unique identifier for the underlying blockchain, 
        represented as the hash of the chain's genesis block encoded in \gls{le} 
        format\footnote{For instance, the \code{chain\_hash} for the Bitcoin Mainnet is 
        \code{6fe28c0ab6f1b372c1a6a246ae63f74f931e8365e15a089c68d6190000000000}~\cite{lightningdevelopers_lightning_2018}.}. 
        This prevents channels established on testnets or forks from being maliciously announced on the main network.

    \item[\code{short\_channel\_id}] 
        The \qty{8}{\byte} base-layer tuple (\term{block height}, \term{transaction index}, and \term{output index}) 
        mapping exactly to the channel's funding transaction.

    \item[\code{node\_id}] 
        The \qty{33}{\byte} compressed public keys of the two participating Lightning nodes. 
        To ensure a deterministic message format and prevent peers from broadcasting duplicate announcements by swapping the fields, 
        \code{node\_id\_1} must be lexicographically strictly less than \code{node\_id\_2}.  

    \item[Bitcoin Keys] 
        The \qty{33}{\byte} compressed public keys utilized in the 2-of-2 multisignature transaction. 
        These are distinct from the Lightning \code{node\_id} keys and are required to verify the on-chain Bitcoin signatures.
\end{description}


\subsubsection{Node Announcement \code{257}}

The \code{node\_announcement} message allows a node to advertise its presence, metadata, and connectivity details to the \gls{ln}.
The protocol specifies the byte layout shown in \autoref{fig:node_announcement}.

\begin{figure}[htbp]
    \centering
    % 32-bit width
    \begin{bytefield}[bitwidth=1.1em, endianness=big]{32}
        \bitheader{0, 15, 16, 31} \\
        
        % Type 
        \wordbox{1}{\code{type} (2 bytes, value: 257)} \\
        
        % Signature 
        \wordbox{2}{\code{signature} (64 bytes)} \\
        
        % Features 
        \bitbox{16}{\code{flen} (2 bytes)} & \bitbox{16}{\code{features} (\code{flen} bytes)} \\
        
        % Timestamp 
        \wordbox{1}{\code{timestamp} (4 bytes)} \\
        
        % ID 
        \wordbox{1}{\code{node\_id} (33 bytes)} \\
        
        % Color & Alias 
        \wordbox{1}{\code{rgb\_color} (3 bytes)} \\
        \wordbox{1}{\code{alias} (32 bytes)} \\
        
        % Addresses 
        \wordbox{1}{\code{addrlen} (2 bytes)} \\
        \wordbox{1}{\code{addresses} (\code{addrlen} bytes)}
    \end{bytefield}
    \caption{
        Byte layout of the \code{node\_announcement} message. 
        Large fields are shown as blocks~\cite[Section~node\_announcement]{lightningdevelopers_lightning_2018}.
    }
    \label{fig:node_announcement}
\end{figure}

The fields are defined as follows:
\begin{description}
    \item[\code{signature}] 
        A \qty{64}{\byte} cryptographic signature over the message hash. 
        This proves that the message originated from the exact \code{node\_id} being announced, 
        preventing peers from spoofing a node's routing configuration.
    
    \item[\code{features}] 
        A \gls{varint} bitfield (dictated by \code{flen}) used to signal supported protocol extensions. 
        As with \code{channel\_announcements}, this ensures forward compatibility across the network.
    
    \item[\code{timestamp}] 
        A \qty{4}{\byte} Unix timestamp. 
        Because nodes dynamically update their network addresses and configurations, 
        this timestamp allows peers to strictly order incoming updates and safely discard stale announcements.

    \item[\code{node\_id}] 
        The \qty{33}{\byte} compressed public key that cryptographically identifies the announcing node on the \gls{ln}.
    
    \item[\code{alias} and \code{rgb\_color}] 
        Metadata allowing operators to customize their node's display name (\tech{UTF-8} encoded) and visual appearance in network explorers. 
        Because these values are entirely arbitrary and unauthenticated, they offer no security guarantees and could be used for impersonation attacks.
    
    \item[\code{addresses}] 
        A \gls{varint} field (dictated by \code{addrlen}) specifying the network transport protocols available for connection. 
        A node may advertise multiple addresses simultaneously, supporting 
        IPv4, IPv6, Tor (v2 and v3) onion services, and DNS hostnames~\cite[Section~node\_announcement]{lightningdevelopers_lightning_2018}.
\end{description}



\subsubsection{Channel Update \code{258}}

Because \gls{ln} channels are bidirectional, routing policies are managed independently for each direction. 
The direction routing from \code{node\_id\_1} to \code{node\_id\_2} is designated as direction \code{0}, 
while the reverse is designated as direction \code{1}. 
After a channel is established, each participant independently broadcasts a \code{channel\_update} message 
to specify their respective fee structure and operational limits~\cite[Section~channel\_update]{lightningdevelopers_lightning_2018}.

\begin{figure}[htbp]
    \centering
    % 32-bit width (4 bytes per row)
    \begin{bytefield}[bitwidth=1.1em, endianness=big]{32}
        
        \bitheader{0,15,16,31} \\
        
        % Type 
        \wordbox{1}{\code{type} (2 bytes, value: 258)} \\
        
        % Signature 
        \wordbox{2}{\code{signature} (64 bytes)} \\
        
        % Chain Hash 
        \wordbox{1}{\code{chain\_hash} (32 bytes)} \\
        
        % SCID (8 bytes = 2 rows) 
        \wordbox{2}{\code{short\_channel\_id} (8 bytes)} \\
        
        % Timestamp 
        \wordbox{1}{\code{timestamp} (4 bytes)} \\
        
        % Flags (1 + 1 + 2 = 4 bytes) 
        \bitbox{8}{\code{msg\_flags} (1 B)} & \bitbox{8}{\code{ch\_flags} (1 B)} & \bitbox{16}{\code{cltv\_expiry\_delta} (2 B)} \\
        
        % Values (u64 = 8 bytes = 2 rows) 
        \wordbox{2}{\code{htlc\_minimum\_msat} (8 bytes)} \\
        \wordbox{1}{\code{fee\_base\_msat} (4 bytes)} \\
        \wordbox{1}{\code{fee\_prop\_millionths} (4 bytes)} \\
        \wordbox{2}{\code{htlc\_maximum\_msat} (8 bytes)}
        
    \end{bytefield}
    \caption{
        Byte layout of the \code{channel\_update} message. 
        Large fields are shown as blocks~\cite[Section~channel\_update]{lightningdevelopers_lightning_2018}.
    }
    \label{fig:channel_update}
\end{figure}

\begin{description}
    \item[Signature] 
        One \qty{64}{\byte} signatures prove that the message has been sent by the nodes that 
        participate in the channel that the update belongs to.

    \item[Metadata] 
        Contains the \code{timestamp} to order updates, the \code{chain\_hash} indicates the correct network and 
        the \code{\gls{scid}} identifies the channel to update.

    \item[Flags] 
        Bitfields containing control information. 
        The most important is the \code{direction} bit (in \code{channel\_flags}), 
        which signals if the update applies to \code{node\_1} $\rightarrow$ \code{node\_2} or the reverse. 
        The \code{disable} bit can be used to temporarily disable the channel.
    
    \item[Fee Policy] 
        Consists of two parameters: \code{fee\_base\_msat} and \code{fee\_proportional\_millionths}. 
        A payment with an amount of $n$\,\unit{\milli\sat} needs to pay a fee following the calculation of \autoref{eq:fee_policy}.
    \begin{align}
        Fee(n) = \code{fee\_base\_msat} + \left( n \cdot \frac{\code{fee\_prop\_millionths}}{1\,000\,000} \right)
        \label{eq:fee_policy}
    \end{align}

    \item[\gls{htlc} Limits] 
        The \code{cltv\_expiry\_delta} defines the required time-lock buffer. 
        The \code{htlc\_minimum\_msat} and \code{htlc\_maximum\_msat} fields define 
        the smallest and largest payment sizes allowed to pass through this channel, respectively.
\end{description}


\subsection{Gossip Synchronization and Queries}
\label{sec:gossip-queries}

While the announcement and update messages described above contain the actual topological data of the \gls{ln}, 
nodes require a standardized mechanism to request this data. 
When a node bootstraps or reconnects after being offline, 
it must synchronize its local graph with the network to ensure its routing tables are accurate. 

To facilitate this, BOLT\#7 defines a suite of query messages that allow nodes to request 
targeted subsets of the gossip data rather than downloading 
the entire network history unnecessarily~\cite[Section~Query Messages]{lightningdevelopers_lightning_2018}.
We divide the query messages into their two-step synchronization workflow:

\begin{enumerate}
    \item \textbf{Macro-Synchronization (\code{query\_channel\_range}):} 
        A node first asks a peer for all active channels created within a specific window of Bitcoin blocks. 
        The peer responds with a \code{reply\_channel\_range} message, returning a compressed array of \glspl{scid} 
        that match the requested timeframe.
    \item \textbf{Micro-Synchronization (\code{query\_short\_channel\_ids}):} 
        Once the requesting node has the list of \glspl{scid}, it compares them against its local database. 
        For any missing or outdated channels, it sends a \code{query\_short\_channel\_ids} message 
        to explicitly request the exact \code{channel\_announcement} and \code{channel\_update} payloads for those specific channels.
\end{enumerate}


\section{Bitcoin Lightning Implementations}
\label{sec:implementations}

There are several (commercial) companies as well as organizations that have 
implemented the described specifications, each with their own technology and different focus,
but all have their codebase open-sourced.
As the following \autoref{tab:lightning-nodes-implementation} shows, there are, as of writing, 
four working implementations and one getting build.

\begin{table}[htpb]
    \centering
    \caption{Comparison of major Bitcoin Lightning Network implementations.}
    \label{tab:lightning-nodes-implementation}
    \renewcommand{\arraystretch}{1.4}
    \begin{tabular}{@{} l l l c @{}}
        \toprule
        \textbf{Implementation} & 
        \begin{tabular}{@{}l@{}}\textbf{Primary} \\ \textbf{Language}\end{tabular} & 
        \begin{tabular}{@{}l@{}}\textbf{Lead} \\ \textbf{Maintainer}\end{tabular} & 
        \textbf{Est.} \\
        \midrule
        \textbf{\gls{lnd}} & Go & Lightning Labs & 2016 \\
        \textbf{Core Lightning (CLN)} & C & Blockstream & 2015 \\
        \textbf{Eclair} & Scala & ACINQ & 2015 \\
        \textbf{\gls{ldk}} & Rust & Spiral & 2019 \\
        \textbf{NLightning} & C\# .NET & N\'{i}ckolas Goline & 2024 \\
        \bottomrule
    \end{tabular}
\end{table}

\paragraph{\gls{lnd}}
Developed by Lightning Labs in Go, \gls{lnd}~\cite{lnddeveloper_lightningnetwork_2026} 
is currently the most widely deployed node implementation on the network. 
Its design focuses on enterprise tooling, featuring a developer-friendly architecture with extensive gRPC and REST APIs.

\paragraph{\gls{cln}}
Maintained primarily by Blockstream, \gls{cln}~\cite{corelightningdevelopers_elementsproject_2026} 
is written in C and follows a UNIX-style philosophy. 
It separates distinct node functions into isolated daemons, 
making it lightweight with a flexible, multi-language plugin system.

\paragraph{Eclair}
Built by ACINQ in Scala, Eclair~\cite{eclairdevelopers_acinq_2026} 
is an enterprise-grade node engineered specifically to handle high-volume routing. 
It is optimized for native mobile integrations and serves as the backbone for 
ACINQ's Phoenix wallet, which shows its capabilities in managing asynchronous connectivity for light clients.

\paragraph{\gls{ldk}}
Unlike standalone node daemons, the \gls{ldk}~\cite{ldkdevelopers_lightningdevkit_2026}, 
managed by Spiral, is a modular, Rust-based development toolkit. 
It gives engineers a flexible API to embed custom Lightning capabilities directly 
into existing applications and wallets, removing the operational overhead of managing a full, localized routing node.

\paragraph{NLightning}
is the newest addition to the ecosystem. NLightning~\cite{goline_ngoline_2026} 
is an actively developing C\# .NET implementation of the BOLT specifications. 
It aims to provide a robust, enterprise-ready implementation designed for 
critical usage within the broader .NET developer ecosystem.


While all implementations successfully execute the BOLT specifications, 
their different architectures make them suitable for different tasks. 
For the purpose of raw gossip data collection, the \gls{cln} implementation presents many advantages. 
Its UNIX-style separation of daemons and its extensible plugin architecture allows 
hooking directly into the node's internal event streams without needing to modify any software. 
Because this design enables real-time access to raw \gls{p2p} gossip messages, 
\gls{cln} serves as the implementation of choice for the \tech{ln-history} platform's data collection.
This will be detailed in the following chapter (\autoref{ch:architecture}).





\iffalse
It is further necessary to specify the 
As soon as a transaction is incorporated into a valid Bitcoin block, meaning a miner has broadcasted a block with a valid block-hash, 
the transaction changes its previously \gls{utxo} to so-called \textit{Spent Transaction Outputs}. Each transaction output represents an amount of bitcoin.
Transactions in their "raw form" can be displayed in hex as seen in \autoref{lst:first-transaction-hex}. 
To better understand this important concept we analyse the very first bitcoin transaction (that was not a coinbase transaction) which occured in block 170 where Satoshi Nakamoto transferred $\si{10}{BTC}$ to Hal Finney.
\begin{listing}[H]
\begin{minted}[
    breaklines=true, 
    breakanywhere=true,
    frame=single, 
    bgcolor=black!5,
    fontsize=\small
]{text}
> bitcoin-cli getrawtransaction f4184fc596403b9d638783cf57adfe4c75c605f6356fbc91338530e9831e9e16

0100000001c997a5e56e104102fa209c6a852dd90660a20b2d9c352423edce25857fcd3704000000004847304402204e45e16932b8af514961a1d3a1a25fdf3f4f7732e9d624c6c61548ab5fb8cd410220181522ec8eca07de4860a4acdd12909d831cc56cbbac4622082221a8768d1d0901ffffffff0200ca9a3b00000000434104ae1a62fe09c5f51b13905f07f06b99a2f7159b2225f374cd378d71302fa28414e7aab37397f554a7df5f142c21c1b7303b8a0626f1baded5c72a704f7e6cd84cac00286bee0000000043410411db93e1dcdb8a016b49840f8c53bc1eb68a382e97b1482ecad7b148a6909a5cb2e0eaddfb84ccf9744464f82e160bfa9b8b64f9d4c03f999b8643f656b412a3ac00000000
\end{minted}
\caption{Retrieving the raw transaction data by txid via bitcoin-cli.}
\label{lst:first-transaction-hex}
\end{listing}

Following the bitcoin protocol this long string of byte data can be organized into 'standardized' fields.
There are three different types of field.
The fixed length, variable length and. 
Fixed length means that there is a pre-set number of bytes allowed for that field .
Variable Length.

The first $\si{4}{B}$ are always the version. Parsing of 01000000 needs to be done using little-endian yielding version: 1.
Next the number of inputs is one byte $01$ meaning 1.
Next $\si{32}{B}$ is the transactionID (c997a5e56e104102fa209c6a852dd90660a20b2d9c352423edce25857fcd3704) always fixed which must be parsed as little endian yielding to: 0437cd7f8525ceed2324359c2d0ba26006d92d856a9c20fa0241106ee5a597c9
This is followed by the output number which is always a fixed 8-byte: 0000000048473044 which cields to 
()

\begin{table}[htpb]
    \centering
    \caption{Hexadecimal breakdown of the first transaction from Satoshi to Hal Finney.}
    \label{tab:tx-hex-breakdown}
    \renewcommand{\arraystretch}{1.3}
    \begin{tabularx}{\textwidth}{|>{\ttfamily\footnotesize\raggedright\arraybackslash}p{3.5cm}|>{\raggedright\arraybackslash}X|l|l|}
        \hline
        \rowcolor{black!10}
        \textbf{Raw Hex Segment} & \textbf{Field Description} & \textbf{Data Type} & \textbf{Size} \\ \hline
        01000000 & \textbf{Version:} 1 & Fixed & 4 Bytes \\ \hline
        01 & \textbf{Input Count:} 1 & CompactSize & 1 Byte \\ \hline
        \multicolumn{4}{|c|}{\cellcolor{blue!10}\textbf{Input 0 (from previous Satoshi transaction)}} \\ \hline
        c997a5e5...7fcd3704 & \textbf{Previous TXID:} (Internal little-endian order) & Fixed & 32 Bytes \\ \hline
        00000000 & \textbf{Previous Output Index:} 0 & Fixed & 4 Bytes \\ \hline
        48 & \textbf{Script Length:} 72 Bytes & CompactSize & 1 Byte \\ \hline
        47304402...1d0901 & \textbf{ScriptSig:} Unlocking script (Signature) & Variable & 72 Bytes \\ \hline
        ffffffff & \textbf{Sequence:} Standard final sequence & Fixed & 4 Bytes \\ \hline
        02 & \textbf{Output Count:} 2 & CompactSize & 1 Byte \\ \hline
        \multicolumn{4}{|c|}{\cellcolor{green!10}\textbf{Output 0 (10 BTC Transferred to Hal Finney)}} \\ \hline
        00ca9a3b00000000 & \textbf{Amount:} $10 \text{ BTC} = 10^9 \text{ sats}$ ($3b9aca00_{16}$) & Fixed & 8 Bytes \\ \hline
        43 & \textbf{Script Length:} 67 Bytes & CompactSize & 1 Byte \\ \hline
        4104ae1a...d84cac & \textbf{ScriptPubKey:} P2PK Locking script & Variable & 67 Bytes \\ \hline
        \multicolumn{4}{|c|}{\cellcolor{red!10}\textbf{Output 1 (40 BTC Change returned to Satoshi)}} \\ \hline
        00286bee00000000 & \textbf{Amount:} $40 \text{ BTC} = 4 \times 10^9 \text{ sats}$ ($ee6b2800_{16}$) & Fixed & 8 Bytes \\ \hline
        43 & \textbf{Script Length:} 67 Bytes & CompactSize & 1 Byte \\ \hline
        410411db...12a3ac & \textbf{ScriptPubKey:} P2PK Locking script & Variable & 67 Bytes \\ \hline
        00000000 & \textbf{Locktime:} 0 & Fixed & 4 Bytes \\ \hline
    \end{tabularx}
\end{table}


\subsubsection{Query and Reply Messages}

A node needs to request its peers to receive new gossip messages. Likewise a peer needs to respond to those \term{Query Messages} using their adequate reply. In \code{BOLT 07} there are three \term{request} and two \term{reply messages} specified and shown in \autoref{tab:gossip-query-reply-bolt-07}.

\begin{table}[htbp]
    \centering
    
    % Define 4 columns: Request Type, Request Name, Reply Type, Reply Name
    \begin{tabular}{@{} ll ll @{}}
        \toprule
        \multicolumn{2}{c}{\textbf{Request}} & \multicolumn{2}{c}{\textbf{Reply}} \\
        \cmidrule(r){1-2} \cmidrule(l){3-4} 
        \textbf{Type} & \textbf{Name} & \textbf{Type} & \textbf{Name} \\
        \midrule
        
        \code{261} & \code{query\_short\_channel\_ids} & \code{262} & \code{reply\_short\_channel\_ids\_end} \\
        \code{263} & \code{request\_channel\_range}    & \code{264} & \code{reply\_channel\_range} \\
        \code{265} & \code{gossip\_timestamp\_filter}  &           & --- \\
        
        \bottomrule
    \end{tabular}

    \caption{Gossip query and reply message types defined in BOLT 07. Note that \code{gossip\_timestamp\_filter} does not trigger a direct reply~\cite[section ~Query Messages]{lightningdevelopers_lightning_2018}.}
    \label{tab:gossip-query-reply-bolt-07}
    
\end{table}

\paragraph{Query Short Channel Ids \code{261}}
~\cite[section ~query\_short\_channel\_ids]{lightningdevelopers_lightning_2018}
\begin{figure}[htbp]
    \centering
    % 32-bit width (4 bytes per row)
    \begin{bytefield}[
        bitwidth=1.1em, endianness=big]{32}
        \bitheader{0,15,16,31} \\
        
        % --- Type ---
        \bitbox{16}{\code{type} (261)} & \bitbox{16}{} \\
        
        % --- Chain Hash (32 bytes) ---
        \wordbox{1}{\code{chain\_hash} (32 bytes)} \\
        
        % --- Len (u16) + Padding for visual alignment ---
        % 'len' is 2 bytes. We place it on the left.
        \bitbox{16}{\code{len} (u16)} & \bitbox{16}{\textit{(reserved space)}} \\
        
        % --- Encoded Short IDs (Variable Length) ---
        % The length is determined by 'len' above
        \wordbox[lrt]{1}{\code{encoded\_short\_ids}} \\
        \wordbox[lrb]{1}{(\code{len} bytes)} \\
        
        % --- TLVs (Variable Stream) ---
        % Contains query_flags (type 1)
        \wordbox[lrt]{1}{\code{tlvs}} \\
        \wordbox[lrb]{1}{(variable length stream)} \\
        
    \end{bytefield}
    \caption{Byte layout of the \code{query\_short\_channel\_ids} message. The \code{encoded\_short\_ids} field contains compressed \gls{scid}s, and the \code{tlvs} field may contain \code{query\_flags} to filter the response~\cite[section ~query\_short\_channel\_ids]{lightningdevelopers_lightning_2018}.}
    \label{fig:query_short_channel_ids}
\end{figure}


\paragraph{Reply Short Channel Ids End \code{262}}

\begin{figure}[htbp]
    \centering
    % 32-bit width (4 bytes per row)
    \begin{bytefield}[
        bitwidth=1.1em, endianness=big]{32}
        \bitheader{0,15,16,31} \\
        
        % --- Type (2 bytes) ---
        \bitbox{16}{\code{type} (262)} & \bitbox{16}{} \\
        
        % --- Chain Hash (32 bytes) ---
        \wordbox{1}{\code{chain\_hash} (32 bytes)} \\
        
        % --- Full Information (1 byte) ---
        % We explicitly show the 1 byte and leave 3 bytes empty for alignment.
        \bitbox{8}{\code{full\_info}} & \bitbox{24}{} \\
        
    \end{bytefield}
    \caption{Byte layout of the \code{reply\_short\_channel\_ids\_end} message. The \code{full\_information} byte (previously named \code{complete}) signals if the replying node believes it has provided all available data~\cite[section ~reply\_short\_channel\_ids\_end]{lightningdevelopers_lightning_2018}.}
    \label{fig:reply_short_channel_ids_end}
\end{figure}


\paragraph{Request Channel Range \code{263}}
\begin{figure}[htbp]
    \centering
    % 32-bit width (4 bytes per row)
    \begin{bytefield}[
        bitwidth=1.1em, endianness=big]{32}
        \bitheader{0,15,16,31} \\
        
        % --- Type ---
        \bitbox{16}{\code{type} (263)} & \bitbox{16}{} \\
        
        % --- Chain Hash ---
        \wordbox{1}{\code{chain\_hash} (32 bytes)} \\
        
        % --- First Blocknum (u32) ---
        \bitbox{32}{\code{first\_blocknum} (u32)} \\
        
        % --- Number of Blocks (u32) ---
        \bitbox{32}{\code{number\_of\_blocks} (u32)} \\
        
        % --- TLVs (Variable) ---
        % Contains query_option (type 1)
        \wordbox[lrt]{1}{\code{tlvs}} \\
        \wordbox[lrb]{1}{(variable stream)} \\
        
    \end{bytefield}
    \caption{Byte layout of the \code{query\_channel\_range} message. The \code{tlvs} field may contain \code{query\_option} flags to request timestamps or checksums for the specified block range~\cite[section ~query\_channel\_range]{lightningdevelopers_lightning_2018}.}
    \label{fig:query_channel_range}
\end{figure}


\paragraph{Reply Chanel Range \code{264}}

\begin{figure}[htbp]
    \centering
    % 32-bit width (4 bytes per row)
    \begin{bytefield}[
        bitwidth=1.1em, endianness=big]{32}
        \bitheader{0,15,16,31} \\
        
        % --- Type ---
        \bitbox{16}{\code{type} (264)} & \bitbox{16}{} \\
        
        % --- Chain Hash ---
        \wordbox{1}{\code{chain\_hash} (32 bytes)} \\
        
        % --- Block Ranges ---
        \bitbox{32}{\code{first\_blocknum} (u32)} \\
        \bitbox{32}{\code{number\_of\_blocks} (u32)} \\
        
        % --- Sync Complete (1 byte) & Len (2 bytes) ---
        % Total 24 bits used. We leave 8 bits empty for visual alignment 
        % before starting the large data blocks.
        \bitbox{8}{\code{sync}} & 
        \bitbox{16}{\code{len} (u16)} & 
        \bitbox{8}{} \\
        
        % --- Encoded Short IDs ---
        \wordbox[lrt]{1}{\code{encoded\_short\_ids}} \\
        \wordbox[lrb]{1}{(\code{len} bytes)} \\
        
        % --- TLVs ---
        % May contain timestamps (type 1) and checksums (type 3)
        \wordbox[lrt]{1}{\code{tlvs}} \\
        \wordbox[lrb]{1}{(timestamps / checksums)} \\
        
    \end{bytefield}
    \caption{Byte layout of the \code{reply\_channel\_range} message. The \code{sync} (\code{sync\_complete}) byte indicates if this is the final reply in a sequence. The \code{tlvs} stream may contain \code{encoded\_timestamps} and \code{checksums} to facilitate efficient reconciliation~\cite[section ~reply\_channel\_range]{lightningdevelopers_lightning_2018}.}
    \label{fig:reply_channel_range}
\end{figure}


\paragraph{Gossip Timestamp Filter \code{265}}
\begin{figure}[htbp]
    \centering
    % 32-bit width (4 bytes per row)
    \begin{bytefield}[
        bitwidth=1.1em, endianness=big]{32}
        \bitheader{0,15,16,31} \\
        
        % --- Type ---
        \bitbox{16}{\code{type} (265)} & \bitbox{16}{} \\
        
        % --- Chain Hash ---
        \wordbox{1}{\code{chain\_hash} (32 bytes)} \\
        
        % --- First Timestamp (u32) ---
        \bitbox{32}{\code{first\_timestamp} (u32)} \\
        
        % --- Timestamp Range (u32) ---
        \bitbox{32}{\code{timestamp\_range} (u32)} \\
        
    \end{bytefield}
    \caption{Byte layout of the \code{gossip\_timestamp\_filter} message. This message constrains relayed gossip to a specific time window defined by \code{first\_timestamp} and the duration \code{timestamp\_range}~\cite[section ~gossip\_timestamp\_filter]{lightningdevelopers_lightning_2018}.}
    \label{fig:gossip_timestamp_filter}
\end{figure}

\fi